% =====================================================================
% Appendix A.2 — Linearity of f1 in drive amplitude
%
% Regenerated by: audit/scripts/appendix_A2_linearity.py  (exit 0)
% Source code reused: four_floor/linearity_check.py (analyse_one, and the
%   pooled-scatter and resolution arithmetic of its has_repeats branch)
% Audit item: 1.8 (promoted from 1.5b)
%
% PACKAGES REQUIRED: booktabs
% OPTIONAL: siunitx (not used below; all units are written out, so this
%   block compiles with booktabs alone)
% PREAMBLE CLASH RISK: none. No new commands are defined. Labels are all
%   prefixed app:lin / tab:lin to avoid collision.
% =====================================================================

\section{Linearity of the first mode in drive amplitude}
\label{app:lin}

Section~4.1.5 reports that the swept-sine replicates show no detected dependence
of $f_1$ on excitation level, and states a resolution of roughly 3.6\%. Both the
test and that resolution figure are set out here, because the resolution is what
the argument in Section~4.3 rests on rather than the $p$-value.

\subsection*{Captures}

Three drive amplitudes were recorded, three replicate sweeps at each, nine
captures in total. Amplitude is reported as the AC RMS of the response, which is
the available proxy for how hard the frame was driven. The ratio of the highest
to the lowest group mean is $2.22$, the factor of 2.2 quoted in Section~4.1.5.

\begin{table}[htbp]
\centering
\caption{Per-capture first-mode estimate and drive level for the nine swept-sine
replicates. Extraction is the project's standard peak-picking over the full band,
$0.5$ to $450$~Hz.}
\label{tab:lin-percapture}
\begin{tabular}{llrr}
\toprule
Gain & Replicate & AC RMS (g) & $f_1$ (Hz) \\
\midrule
1v4 & r1 & 0.32234 & 2.882 \\
1v4 & r2 & 0.32420 & 2.819 \\
1v4 & r3 & 0.32995 & 2.950 \\
\midrule
2v2 & r1 & 0.46840 & 3.044 \\
2v2 & r2 & 0.48162 & 3.044 \\
2v2 & r3 & 0.47505 & 2.962 \\
\midrule
2v8 & r1 & 0.74291 & 2.879 \\
2v8 & r2 & 0.72004 & 2.877 \\
2v8 & r3 & 0.70295 & 2.812 \\
\bottomrule
\end{tabular}
\end{table}

Note that the drive level is nearly constant within a gain group, spanning at
most $3\%$, while $f_1$ varies by up to $0.13$~Hz within the same group. The
within-group variation is not a drive effect.

\subsection*{Group statistics}

\begin{table}[htbp]
\centering
\caption{First-mode estimate by drive level. The standard deviation is over the
three replicates at each gain and measures the repeatability of the swept-sine
estimate at fixed drive.}
\label{tab:lin-groups}
\begin{tabular}{lrrrrr}
\toprule
Gain & $n$ & Mean RMS (g) & Mean $f_1$ (Hz) & SD $f_1$ (Hz) & Range (Hz) \\
\midrule
1v4 & 3 & 0.32549 & 2.884 & 0.0658 & 0.1316 \\
2v2 & 3 & 0.47502 & 3.017 & 0.0475 & 0.0826 \\
2v8 & 3 & 0.72197 & 2.856 & 0.0381 & 0.0668 \\
\bottomrule
\end{tabular}
\end{table}

\subsection*{Tests}

Two tests are reported, because they disagree and the disagreement is
informative.

The test quoted in Section~4.1.5 compares the extreme gains only, by Welch's
two-sample $t$-test: $t = -0.63$, $\mathrm{df} = 3.21$, $p = 0.573$. It discards
the middle gain.

A one-way analysis of variance across all three gains gives
$F(2,6) = 8.24$, $p = 0.019$. The between-group differences are therefore
significant, and the extreme-gain test cannot see them: the 1v4 and 2v8 means
differ by only $0.028$~Hz, while the 2v2 mean sits $0.147$~Hz above their
midpoint. The deviating group is the middle one.

That significant variation is not an amplitude effect. A linear regression of
$f_1$ on drive RMS over all nine captures gives a slope of $-0.120$~Hz per unit
RMS with standard error $0.183$, $t = -0.66$ on 7 degrees of freedom, $p = 0.533$,
$r = -0.24$. The group means are not monotone in drive. The consistent reading is
that the swept-sine estimate of $f_1$ carries a run-to-run component larger than
its within-run scatter, which is a property of the estimator rather than of the
structure.

\subsection*{Resolution}

The pooled within-gain standard deviation is $0.0518$~Hz on 6 degrees of freedom.
Against the lowest-gain mean of $2.8836$~Hz that is $1.79\%$ of $f_1$, and twice
it is $3.59\%$, the figure quoted in Section~4.1.5 and in the assumptions table.
The derivation is therefore
%
\begin{equation*}
  \text{resolution} \;=\; \frac{2\,s_{\text{pooled}}}{\bar{f}_{1,\text{lowest gain}}}
  \;=\; \frac{2 \times 0.0518}{2.8836} \;=\; 3.59\%,
\end{equation*}
%
a two-standard-deviation figure, normalised by the lowest-gain group mean.

Two standard deviations is a rule of thumb rather than a formal detection limit.
The minimum difference this design could detect in a two-sample comparison at
three replicates per group, $\alpha = 0.05$ and $80\%$ power, is $0.153$~Hz, or
$5.30\%$ of $f_1$. The formal limit is therefore weaker than the quoted $3.59\%$,
which means the test is less sensitive than that figure suggests, not more.

\subsection*{What controls the trace grade}

The resolution above is far coarser than three of the four trace-grade shifts, so
it cannot be what licenses them. Those rest instead on the repeatability of the
tap-based estimate within a cell, which is measured directly and is an order of
magnitude tighter.

\begin{table}[htbp]
\centering
\caption{Within-cell tap scatter of $f_1$ at the four trace-grade cells, over the
five taps of each cell. Compare the $3.59\%$ resolution of the swept-sine test.}
\label{tab:lin-trace}
\begin{tabular}{lrrr}
\toprule
Cell & $n$ & Mean $f_1$ (Hz) & Tap SD (\% of mean) \\
\midrule
Base plate, trace & 5 & 2.8548 & 0.258 \\
Floor 1, trace    & 5 & 2.7968 & 0.054 \\
Floor 2, trace    & 5 & 2.8932 & 0.171 \\
Floor 3, trace    & 5 & 2.9208 & 0.226 \\
\bottomrule
\end{tabular}
\end{table}

\subsection*{Conclusion}

No dependence of $f_1$ on drive amplitude was detected across a factor of 2.2 in
drive. This is an absence of a detected effect and not a demonstration of
linearity: the design resolves shifts above roughly $3.6\%$ by a two-sigma
criterion, and above roughly $5.3\%$ formally, so an amplitude dependence smaller
than that would not have been found. The graded-series conclusions of Chapter~4
do not rest on this test. Tap amplitude across the graded series spans a factor
of $2.5$ by set mean, inside the range tested here, and the within-cell tap
scatter of Table~\ref{tab:lin-trace} bounds the combined effect of drive
variation and estimator noise at below $0.3\%$ at every trace cell.
